% !TeX program = pdflatex
% ─────────────────────────────────────────────────────────────────────────────
%  lcs_dobra.tex — Longest Common Subsequence pela teoria da dobra
%
%  Compila: pdflatex lcs_dobra.tex
% ─────────────────────────────────────────────────────────────────────────────
\documentclass[11pt,a4paper]{article}
\usepackage[utf8]{inputenc}\usepackage[T1]{fontenc}\usepackage[brazil]{babel}
\usepackage{amsmath,amssymb,amsthm,amscd}
\usepackage[margin=2.4cm]{geometry}
\usepackage{booktabs}\usepackage{xcolor}
\usepackage[hidelinks]{hyperref}

\definecolor{cinza}{gray}{0.35}
\newcommand{\code}[1]{\texttt{\small #1}}
\newcommand{\abs}[1]{\lvert #1\rvert}
\newcommand{\N}{\ensuremath{\mathbb{N}}}
\newtheorem{teorema}{Teorema}
\newtheorem{lema}[teorema]{Lema}
\newtheorem{definicao}[teorema]{Definição}
\newtheorem{corolario}[teorema]{Corolário}
\newtheorem{proposicao}[teorema]{Proposição}
\renewcommand{\proofname}{Demonstração}

\begin{document}

\title{Longest Common Subsequence\\[0.4em]
\large pela teoria da dobra e da multiplicidade geométrica}
\author{Núcleo a partir de \texttt{aranha.tex} / \texttt{arquitetura.tex}}
\date{}
\maketitle

\begin{abstract}
\noindent Duas sequências finitas são duas \textbf{realizações}
$\pi_A\colon I_A\to\Sigma$ e $\pi_B\colon I_B\to\Sigma$.
Uma subsequência comum é um par de mapas estritamente crescentes
que se compõem na mesma palavra. O comprimento máximo é o cardinal
da maior cadeia de posições de coincidência que respeitam a ordem
produto. A recorrência clássica de programação dinâmica é reescrita
como \textbf{dobra} alternada das duas faces (avançar em $A$, avançar
em $B$, ou avançar nas duas quando há match). Porque a régua é finita
o processo \textbf{termina} numa igualdade; o valor $L(m,n)$ é o
ponto fixo comum às duas faces no canto $(m,n)$ da grade.
A recuperação da sequência é o levantamento da fibra (coordenada de
folha) ao longo do caminho óptimo.
\end{abstract}

% ═══════════════════════════════════════════════════════════════════════════
\section{Objectos}

\begin{definicao}[sequência como realização]
Seja $\Sigma$ um alfabeto finito com igualdade decidível.
Uma sequência $A=a_1\dots a_m$ é a realização
\[
\pi_A\colon I_A\to\Sigma,\qquad
I_A=\{1,\dots,m\},\qquad
\pi_A(i)=a_i.
\]
Analogamente $\pi_B\colon I_B\to\Sigma$ para $B=b_1\dots b_n$.
A \textbf{multiplicidade} de um símbolo $\sigma\in\Sigma$ em $A$ é
\[
G_A(\sigma)=\bigl\lvert\pi_A^{-1}(\sigma)\bigr\rvert.
\]
\end{definicao}

\begin{definicao}[subsequência comum]
Uma \textbf{subsequência comum} de comprimento $k$ é um par de mapas
estritamente crescentes
\[
\iota\colon\{1,\dots,k\}\hookrightarrow I_A,\qquad
\kappa\colon\{1,\dots,k\}\hookrightarrow I_B
\]
tais que
\[
\pi_A\circ\iota \;=\; \pi_B\circ\kappa.
\]
O valor comum é a palavra $w=\pi_A(\iota(1))\dots\pi_A(\iota(k))$.
O \textbf{LCS} é qualquer subsequência comum de comprimento máximo;
denotamos esse comprimento por $L(A,B)$.
\end{definicao}

\noindent Em linguagem da aranha: $\iota$ e $\kappa$ são dois
levantamentos parciais; a condição $\pi_A\circ\iota=\pi_B\circ\kappa$
diz que as duas realizações coincidem sobre a imagem comum.

% ═══════════════════════════════════════════════════════════════════════════
\section{A grade e os matches}

\begin{definicao}[poset dos matches]
Seja
\[
M=\bigl\{(i,j)\in I_A\times I_B:\pi_A(i)=\pi_B(j)\bigr\}.
\]
Ordene $M$ pela ordem produto:
\[
(i,j)\;\le\;(i',j')\quad\iff\quad i\le i'\;\text{e}\;j\le j'.
\]
Uma cadeia em $M$ é exactamente uma subsequência comum.
Logo
\[
\boxed{\;L(A,B)\;=\;\text{comprimento da maior cadeia em }(M,\le).\;}
\]
\end{definicao}

\noindent Isto é Dilworth/Erdős–Szekeres no poset dos pontos de
coincidência. A multiplicidade geométrica $G$ da realização
conjunta $\pi(i,j)=(\pi_A(i),\pi_B(j))$ conta quantas vezes cada
par de símbolos aparece; $M$ é o suporte onde a diagonal de $\Sigma$
é atingida.

% ═══════════════════════════════════════════════════════════════════════════
\section{A dobra que calcula o comprimento}

\begin{definicao}[função de estado]
Para $0\le i\le m$, $0\le j\le n$ defina
\[
L(i,j)\;=\;L\bigl(a_1\dots a_i,\;b_1\dots b_j\bigr)
\]
(com $L(0,\cdot)=L(\cdot,0)=0$).
\end{definicao}

\begin{teorema}[recorrência como dobra]\label{thm:lcs-dobra}
Vale a recorrência
\[
L(i,j)\;=\;
\begin{cases}
L(i-1,j-1)+1 & \text{se }\pi_A(i)=\pi_B(j),\\[4pt]
\max\bigl(L(i-1,j),\,L(i,j-1)\bigr) & \text{caso contrário.}
\end{cases}
\]
\end{teorema}

\begin{proof}
Caso match: qualquer cadeia óptima que use o match $(i,j)$ pode ser
obtida de uma cadeia óptima em $M\cap([1..i-1]\times[1..j-1])$
acrescentando esse ponto; e nenhuma cadeia pode usar dois matches
com a mesma primeira ou segunda coordenada. Logo o óptimo com o
match é $L(i-1,j-1)+1$. Sem o match o óptimo é o melhor dos dois
prefixos.

Caso sem match: nenhuma cadeia pode terminar em $(i,j)$, logo o
óptimo é o máximo dos dois prefixos imediatos.
\end{proof}

\medskip\noindent\textbf{Leitura pela teoria da dobra.}
Há três movimentos possíveis na grade:
\begin{center}
\begin{tabular}{@{}lll@{}}
\toprule
movimento & o que faz & face \\
\midrule
$(i-1,j)\to(i,j)$ & avança só em $A$ & face de $A$ \\
$(i,j-1)\to(i,j)$ & avança só em $B$ & face de $B$ \\
$(i-1,j-1)\to(i,j)$ & avança nas duas (match) & ponto fixo comum \\
\bottomrule
\end{tabular}
\end{center}
Quando há match as duas faces coincidem e o valor sobe de $1$
(o ``meio'' produzido pela dobra). Quando não há match toma-se o
máximo --- a face que ainda conserva o maior comprimento.
O processo é exactamente o \textbf{corte} do \texttt{aranha}
(Teor.~corte): batendo as duas faces alternadamente os intervalos
de possíveis comprimentos encaixam e, porque a grade é finita,
\textbf{terminam} numa igualdade $L(m,n)$.

\begin{corolario}[terminação]
O valor $L(m,n)$ é atingido após exactamente $mn$ passos de
preenchimento da tabela (ou $O(mn)$ operações). Não há convergência
assintótica: a régua tem grão e o processo pára.
\end{corolario}

% ═══════════════════════════════════════════════════════════════════════════
\section{Recuperação da sequência --- o levantamento}

Conhecer só $L(m,n)$ dá o comprimento. Para obter uma palavra
óptima é preciso o \textbf{levantamento} (Teor.~levantamento da
aranha): guardar, em cada célula, qual face produziu o valor.

\begin{proposicao}[caminho óptimo]
Defina o ponteiro
\[
P(i,j)\;=\;
\begin{cases}
\leftarrow & \text{se }L(i,j)=L(i,j-1),\\
\uparrow   & \text{se }L(i,j)=L(i-1,j),\\
\nwarrow  & \text{se }L(i,j)=L(i-1,j-1)+1\text{ e match}.
\end{cases}
\]
(Em caso de empate qualquer escolha é válida.)
Partindo de $(m,n)$ e seguindo $P$ até $(0,0)$, cada passo
$\nwarrow$ contribui um símbolo $\pi_A(i)=\pi_B(j)$ para a
subsequência. O caminho é uma cadeia máxima em $M$.
\end{proposicao}

\noindent Isto é o levantamento $\tilde\pi$ que acrescenta a
coordenada de folha ``de onde vim''; a volta é exacta porque cada
ponteiro é uma involução local (desfazer o passo devolve a célula
anterior).

% ═══════════════════════════════════════════════════════════════════════════
\section{Forma matricial e conservação}

Escreva a tabela $L$ como matriz $(m+1)\times(n+1)$ de inteiros
não-negativos. A recorrência é uma operação local
\[
L_{i,j}\;=\;f\bigl(L_{i-1,j},\,L_{i,j-1},\,L_{i-1,j-1},\,\mathbf{1}_{a_i=b_j}\bigr).
\]
Em particular, quando se percorre a tabela por diagonais
$i+j=\mathrm{const}$, cada valor depende só de valores já
computados --- a ordem de avaliação é uma realização da ordem
produto.

\begin{lema}[conservação de diferença]
Para quaisquer $i,j$,
\[
0\;\le\;L(i,j)-L(i-1,j)\;\le\;1,\qquad
0\;\le\;L(i,j)-L(i,j-1)\;\le\;1.
\]
\end{lema}

\begin{proof}
Segue por indução da recorrência: o valor sobe no máximo $1$ e nunca
desce.
\end{proof}

\noindent Logo a ``derivada discreta'' (diferença finita) é $0$ ou
$1$ --- exactamente o bit que a face decide. Isto é a versão
discreta da conservação da norma: a variação total ao longo de
qualquer caminho monótono de $(0,0)$ a $(m,n)$ é $L(m,n)$, e o
número de passos diagonais óptimos é o mesmo.

% ═══════════════════════════════════════════════════════════════════════════
\section{Complexidade e régua}

\begin{proposicao}
Tempo $O(mn)$, espaço $O(\min(m,n))$ (duas linhas bastam, ou uma
com cuidado). O resultado é exacto: não há aproximação nem
truncamento.
\end{proposicao}

\noindent Se a máquina trabalha com $\mathrm{Word}_8$ e a tabela não
cabe, a arquitectura dos papers \textbf{recusa} em vez de truncar
(§ISA). O comprimento $L(m,n)$ cabe sempre num inteiro de
$\lceil\log_2(\min(m,n)+1)\rceil$ bits.

% ═══════════════════════════════════════════════════════════════════════════
\section{Exemplo mínimo}

Sejam $A=\texttt{ABCBDAB}$, $B=\texttt{BDCABA}$.
\[
\begin{array}{c|ccccccc}
 & \varepsilon & B & D & C & A & B & A\\\hline
\varepsilon & 0 & 0 & 0 & 0 & 0 & 0 & 0\\
A & 0 & 0 & 0 & 0 & 1 & 1 & 1\\
B & 0 & 1 & 1 & 1 & 1 & 2 & 2\\
C & 0 & 1 & 1 & 2 & 2 & 2 & 2\\
B & 0 & 1 & 1 & 2 & 2 & 3 & 3\\
D & 0 & 1 & 2 & 2 & 2 & 3 & 3\\
A & 0 & 1 & 2 & 2 & 3 & 3 & 4\\
B & 0 & 1 & 2 & 2 & 3 & 4 & 4\\
\end{array}
\]
$L(7,6)=4$. Um caminho óptimo recupera, por exemplo, $\texttt{BCBA}$
ou $\texttt{BDAB}$.

% ═══════════════════════════════════════════════════════════════════════════
\section{Síntese na linguagem dos papers}

\begin{center}
\begin{tabular}{@{}ll@{}}
\toprule
conceito do paper & no LCS \\
\midrule
realização $\pi$ & as duas sequências \\
multiplicidade $G$ & contagem de cada símbolo \\
poset / ordem produto & matches ordenados \\
maior cadeia & LCS \\
dobra / duas faces & os três movimentos da grade \\
corte que termina & $L(m,n)$ exacto \\
levantamento (folha) & ponteiros $P(i,j)$ \\
conservação & diferenças $0$ ou $1$ \\
régua finita & terminação em $mn$ passos \\
\bottomrule
\end{tabular}
\end{center}

\medskip
\noindent\textbf{Em uma frase:}
o LCS é o comprimento da maior cadeia de coincidências no produto
das duas realizações; a tabela dinâmica é a execução da dobra
alternada das duas faces sobre a grade finita; a sequência em si é
o levantamento dessa cadeia.

\vfill
{\footnotesize\color{cinza}
Referências de estilo e enunciados: \texttt{papers/aranha.tex}
(dobra, multiplicidade, levantamento, corte),
\texttt{papers/arquitetura.tex} (régua, conservação, recusa).
Números de verificação: qualquer implementação clássica de LCS
com os mesmos exemplos.}

\end{document}
